\pagenumbering{gobble}

\selectlanguage{magyar}
\hungarianParagraph

%----------------------------------------------------------------------------
% Abstract in Hungarian
%----------------------------------------------------------------------------

\chapter*{Kivonat}
\begin{spacing}{1}

A mai rohanó világban az ingatlanügyletekben részt vevő magánszemélyek és vállalkozások gyakran jelentős kihívásokkal szembesülnek az ingatlan- és tulajdonadásvételi szerződések elkészítésekor. Ezen szerződések létrehozása az ingatlanügyletek szerves részét képzi, összetett jogi szabályok böngészésével és aprólékos dokumentáció írását köteleztetve. Szükség van egy olyan intuitív eszköz kitalálására, amely enyhíti ezeket a bonyolult folyamatok végigjárását, és megkönnyíti a jogszabályoknak megfelelő, pontos szerződések létrehozását.

Dolgozatom témája kettős. Először is, meg kívánom teremteni egy olyan online platform alapjait, amely lehetővé teszi a szerződések automatikus létrehozását a biztonságos ingatlanértékesítések megkönnyítése érdekében. Ez magában foglalja a szükséges szoftverarchitektúra megtervezését, egy könnyen kezelhető adminisztrációs felület létrehozását, valamint PDF dokumentumok szó szerinti generálását és kitöltését az adminisztrátorok által létrehozott sablonok segítségével. Másodszor, szándékomban áll dokumentálni a végső kitelepítés különböző szakaszait (és lehetőségeit), valós körülmények között mérve.

A javasolt weboldal átfogó platformként szolgál, barátságos felületet kínálva a felhasználóknak a saját igényeikre szabott, jogilag teljes bizonyító erejű szerződések létrehozására. Azáltal, hogy nincs szükség kiterjedt jogi ismeretekre vagy költséges ügyvédek fizetésére, a szerződések létrehozásához szükséges idő és erőfeszítés jelentősen csökken.

Összehasonlításképpen, világszerte hasonló szoftvermegoldásokat vezetett be számos kormány az adásvételi szerződésekkel kapcsolatos kihívások kezelésére. Ezeket a kormányok által támogatott törekvéseket azonban gyakran jelentősen hátráltatják a meggondolatlan tervezési döntéseik és az eredendő homályosságuk. Platformunk célja, hogy a felhasználók számára olyan egyedi leegyszerűsített élményt nyújtson, amely felülmúlja a meglévő kormányzati kezdeményezéseket.

A szakdolgozat elején a javasolt szoftveres megoldás létrehozásához alkalmazott technológiákat tárgyaljuk. Ezt követően leírjuk, hogy mit fog a szoftver szolgáltatni, és hogy milyen teljesítményt várunk el tőle egy szoftverkövetelményspecifikáció részeként. A lehetséges szoftverarchitektúra és telepítési bemutatását követően részletesen bemutatjuk a platform leendő felhasználói felületét. Végezetül bemutatunk néhány mérést a szoftverarchitektúra tervezési szakaszában hozott döntéseink racionalizálása érdekében.

\vskip 0.5cm

\textbf{Kulcsszavak}: \textit{Tulajdonszerződések}, \textit{Dokumentumsablonok}, \textit{PDF dokumentumgenerálás}, \textit{Dockeres kitelepítés}, \textit{Serverless technológiák}.
\end{spacing}
\vfill
\selectlanguage{romanian}

%----------------------------------------------------------------------------
% Abstract in Romanian
%----------------------------------------------------------------------------
\chapter*{Rezumat}
\begin{spacing}{1}

În ziua de azi, persoanele fizice și juridice angajate în tranzacții imobiliare se confruntă adesea cu provocări semnificative atunci când redactează contracte de vânzare de proprietăți. Crearea acestor contracte de vânzare a proprietății este partea integrală a tranzacțiilor imobiliare, implicând legalități complexe și necesitatea unei documentații meticuloasă. Există o nevoie pentru un instrument intuitiv care să atenueze aceste complexități și să faciliteze crearea unor contracte precise și conforme cu legea.

Scopul tezei mele este dublu. În primul rând, intenționez să creez bazele unei platforme online care să permită crearea automată de contracte pentru a facilita vânzarea sigură a proprietăților. Aceasta include planificarea unei arhitecturii software necesare, crearea unei interfețe de administrare ușor de utilizat, precum și generarea și completarea propriu-zisă a documentelor PDF, folosind șabloanele create de către administratori. În al doilea rând, intenționez să documentez diferitele etape (și posibilități) ale instalării finale, măsurate în condiții reale.

Site-ul web propus servește ca o platformă cuprinzătoare, oferind utilizatorilor o interfață prietenoasă pentru a genera contracte obligatorii din punct de vedere juridic, adaptate la cerințele lor specifice. Prin eliminarea necesității unor cunoștințe juridice extinse sau a asistenței unor avocați costisitori, timpul și efortul necesar pentru a crea un contract sunt mult reduse.

Comparativ, guvernele din întreaga lume au implementat soluții software asemănătoare pentru a aborda provocările asociate cu contractele de vânzare de proprietăți. Cu toate acestea, aceste eforturi susținute de guverne sunt adesea dezavantajate în mod semnificativ de obscuritatea lor percepută. Această platformă își propune să ofere utilizatorilor o experiență unică și simplificată, care depășește inițiativele guvernamentale existente.

La începutul tezei vom discuta diferitele tehnologii utilizate pentru a crea soluția software propusă. Ulterior, vom descrie ceea ce va face software-ul și cum se așteaptă să funcționeze ca parte a unei specificații a cerințelor software (\textit{SRS}). O interfață de utilizator prospectivă pentru platformă va fi detaliată în urma propunerii unei arhitecturi software. În cele din urmă, vor fi prezentate câteva măsurători pentru a raționaliza deciziile luate în timpul etapelor de planificare a arhitecturii software.

\vskip 0.5cm

\textbf{Cuvinte cheie}: \textit{Contracte de proprietate}, \textit{Șabloane}, \textit{Generarea de documente PDF}, \textit{Instalare folosind Docker}, \textit{Tehnologii serverless}.
\end{spacing}
\vfill
\selectlanguage{english}
%\englishParagraph

%----------------------------------------------------------------------------
% Abstract in English
%----------------------------------------------------------------------------
\chapter*{Abstract}
\begin{spacing}{1}

In today's fast-paced world, individuals and businesses engaged in real estate transactions often encounter significant challenges when drafting property sale contracts. The creation of these property sale contracts is an integral part of real estate transactions, involving complex legalities and meticulous documentation. There is a pressing need for an intuitive tool that can alleviate these complexities and facilitate the creation of accurate and legally compliant contracts.

The aim of my thesis is two-fold. Firstly, I intend to create the basis of an online platform to allow the automatic creation of contracts to facilitate secure property sales. This includes the creation and planification of the necessary software architecture, the creation of an easy to use administration interface, as well as the verbatim generation and filling out of PDF documents using templates created by the administrators. Secondly, I intend to document the various stages (and possibilities) of the final deployment, measured within real-world conditions.

The proposed website serves as a comprehensive platform, offering users a friendly interface to generate legally binding contracts tailored to their specific requirements. By eliminating the need for extensive legal knowledge or the assistance of costly attorneys, the time and effort required to create a contract is greatly reduced.

Comparatively, governments worldwide have implemented similar software solutions to address the challenges associated with property sale contracts. However, these government-backed endeavours are often significantly handicapped by their perceived obscurity and their convoluted design choices. This platform aims to provide users with a unique and simplified experience that surpasses existing government initiatives.

In the beginning of the thesis we will discuss the various technologies employed to create the proposed software solution. Afterwards, we will describe what the software will do and how it will be expected to perform as part of a Software Requirements Specification. A prospective user interface for the platform will be detailed following the proposal of a potential software architecture and deployment layout. Finally, some measurements will be presented to rationalize the decisions taken during the planning stages of the software architecture.

\vskip 0.5cm

\textbf{Keywords}: \textit{Property contracts}, \textit{Document templates}, \textit{PDF document generation}, \textit{Docker deployment}, \textit{Serverless technologies}.
\end{spacing}
\vfill
\dolgozatnyelve
\defaultParagraph